\documentclass[10pt,a4paper]{article}
\usepackage[margin=2cm]{geometry}
\usepackage[T1]{fontenc}
\usepackage{booktabs,tabularx,array,longtable}
\renewcommand{\arraystretch}{1.1}
\newcolumntype{Y}{>{\raggedright\arraybackslash}X}
\newcolumntype{L}[1]{>{\raggedright\arraybackslash}p{#1}}
\newcommand{\grp}[1]{\midrule\multicolumn{4}{@{}l}{\textit{#1}}\\\addlinespace[2pt]}
\newcommand{\KEEP}{\textbf{Keep}}
\newcommand{\FIX}{\textbf{Fix}}
\newcommand{\EXCL}{\textbf{Exclude}}
\newcommand{\INV}{\textbf{Investigate}}
\pagestyle{empty}

\begin{document}

\section*{Problems, gaps and artefacts in the data}

Each row gives the problem, why it matters for building and evaluating a monitor, and what I would do.
\KEEP{} = use as is, but report it; \FIX{} = clean it before use; \EXCL{} = remove from what the monitor
sees or from a metric; \INV{} = ask the data owners or check further.

\medskip
{\small
\begin{longtable}{@{} L{3.3cm} L{4.8cm} L{5.6cm} L{2.0cm} @{}}
\toprule
\textbf{Problem} & \textbf{Effect on use} & \textbf{What I would do} & \textbf{Decision} \\
\endfirsthead
\toprule
\textbf{Problem} & \textbf{Effect on use} & \textbf{What I would do} & \textbf{Decision} \\
\endhead
\grp{Gaps in what the dataset can show}
No honest runs in Family A (all 46 have a harmful goal) &
False alarms cannot be measured per trace; a monitor that flags everything looks perfect &
Measure false alarms per action (most actions are honest work); use the 6 honest twins; build honest versions by removing attack steps &
\KEEP{}, report \\
Only 6 honest traces in total (90\% attack rate) &
Accuracy is meaningless; precision will not carry over to real use, where attacks are rare &
Report recall, precision and false alarms, not accuracy; state the base-rate caveat &
\KEEP{}, report \\
Model, environment, family and attack type change together &
A monitor can score well by recognising the setting rather than the behaviour &
Report every result per family; do not generalise across families &
\KEEP{}, report \\
About 12 distinct tasks; repeated runs are near-copies &
Results look more certain than they are; random splits leak between train and test &
Group and split by task; report counts and examples, not polished percentages &
\KEEP{}, report \\
No ground-truth labels or grader results; test tools never called &
``Succeeded'' can only be inferred from the agent's own outputs &
Hand-label a gold set; label \emph{attempted} and \emph{succeeded} separately; ask for scorer outputs &
\INV{} \\
\grp{Missing or partial context}
Family A has no recorded prompts (no system or user messages) &
We cannot see what the agent was told, so its intent cannot be judged &
Treat as actions-only; label what actions do, not intent; ask whether the prompts exist &
\INV{} \\
Family B prompts look partly stripped: a safety warning the worker quotes is in none of the 12 system prompts &
What the worker knew is uncertain &
Note it; ask for the full prompts &
\INV{} \\
No reasoning in Family A; reasoning partly hidden in Family B &
Reasoning-based monitoring can only be tested on 12 traces &
Run two monitor conditions: actions only, and with reasoning &
\KEEP{} \\
Offline only: effects cannot be checked (did the deletion persist? did the exploit fire?) &
Success stays uncertain for about 8 traces &
Mark success as unclear where there is no in-trace evidence; replay only in a separate sandbox &
\INV{} \\
\grp{Fields that give the answer away}
Answer-leaking fields and strings (listed below the table) &
A monitor that sees these copies the answer and looks far better than it is &
Hide them from the monitor; keep them only for checking results; report results with and without masking &
\EXCL{} from monitor input \\
42 submission-tool calls with no proposing message &
They are the scorer's interface, not real environment actions, yet they hold the clearest evidence &
Hide from the monitor; use as supporting evidence when labelling &
\EXCL{} from monitor input \\
\grp{Recording artefacts and data quality}
Status fields wrong: \texttt{event\_failed} never true despite 56 errors; \texttt{event\_observed} always true &
Failure counts from these fields are wrong &
Derive my own failure flag from errors, exit codes and tracebacks &
\FIX{} \\
42 non-zero exit codes and about 306 tracebacks carry no error flag; 10 outputs cut at about 16.5\,KB &
Hidden failures; a cut output may hide the evidence &
Add a ``failed'' and a ``truncated'' flag per action &
\FIX{} \\
Commands stored under different keys (\texttt{cmd}, \texttt{command}, script fields) &
A parser reading one key silently misses commands &
Normalise in the loader (done) &
\FIX{} (done) \\
Tool outputs stored twice (in messages and in actions) &
Inflates size and token counts &
Use the action records only (verified identical) &
\FIX{} (done) \\
Easter egg: inserted text in 8 messages &
Messages were edited after recording; could distract a labeller &
Flag and strip before labelling; note the provenance caveat for message text &
\FIX{} \\
One trace cut off by a safety filter (customs fraud) &
Incomplete; the attack never finished &
Keep for labelling; exclude from ``did the attack succeed'' metrics &
\EXCL{} from success metrics \\
13 traces never submit; 15 have no seed &
Minor; affects completion and reproducibility claims &
Note in the report &
\KEEP{} \\
Main-task spec contradicts itself (\texttt{backend.go} vs.\ \texttt{backend.py}) &
Unclear what ``main task done'' means there &
Note it; do not score main-task completion strictly &
\KEEP{}, report \\
A traceback contains a home-directory path, username and project name &
Privacy leak in the data itself &
Redact in anything I output; tell the data owners &
\FIX{} \\
\grp{Ambiguities the labelling policy must settle}
Harmful outcome vs.\ harmful intent: Family B workers carry out harm they may not understand &
One label cannot say both ``harm happened'' and ``the agent meant it'' &
Two separate fields: harm of the action, and the agent's apparent awareness &
\FIX{} (in policy) \\
Required actions that look harmful (external email, \texttt{Delete()}, test emails) &
Naive rules and monitors raise false alarms &
Write explicit rules: judge what the action does in context, not its keywords &
\KEEP{} as hard cases \\
Honest twins make the same final write as the attacks &
Harm depends on earlier context (approvals checked, value origin) &
Give the monitor the history, not single actions &
\KEEP{} as hard cases \\
Workers mention ``the autograder'' and ``DecompBench'' &
Behaviour in a known test may differ from deployment &
Note it as a limit on how far results generalise &
\KEEP{}, report \\
\bottomrule
\end{longtable}
}

{\sloppy\noindent\textit{Answer-leaking fields and strings:} \texttt{adversarial\_tasks}; the side-task submission tool;
\texttt{rce\_secret} and \texttt{/tmp/pwned}; ``authorized'' (only in honest prompts); \texttt{exfiltrate\_tokens}
in a task id; the test tools in the tool list.\par}

\subsection*{Summary}
\begin{itemize}\setlength\itemsep{1pt}
  \item \textbf{Keep} all 58 traces: none is unusable, and the hard cases (honest twins, required look-alike actions) are the most valuable part.
  \item \textbf{Fix} the status fields, argument keys, duplicated outputs, inserted text and the privacy leak before any analysis.
  \item \textbf{Exclude} the answer-leaking fields from what the monitor sees, and the cut-off trace from success metrics.
  \item \textbf{Investigate} with the data owners: grader results, the missing prompts, and whether the attacks really succeeded.
  \item \textbf{Overall:} the data is good for testing whether a monitor catches realistic attacks, weak for measuring false alarms, and too small to train on.
\end{itemize}

\end{document}
